\documentclass[12pt, letterpaper]{article}

%-------Packages---------
\usepackage{amssymb,amsfonts}
\usepackage{mathptmx}
\usepackage{amsthm}
\usepackage{graphicx}
\usepackage[all,arc]{xy}
\usepackage[colorlinks=true, allcolors=blue]{hyperref}
\usepackage{enumerate}
\usepackage{bbm}
\usepackage{dsfont}
\usepackage{mathrsfs}
\usepackage{tikz-cd}
\usepackage{setspace}
\usepackage{import}
\usepackage[utf8]{inputenc}
\usepackage{hyperref}
\usepackage{tikz}
\usepackage{float}
\usepackage{mdframed}
\usepackage{fancyhdr}
\usepackage[nointegrals]{wasysym}

\usepackage[left=1in,top=1in,right=1in,bottom=1in]{geometry}

\usepackage{mathtools}
\usepackage{setspace}
\usepackage{cleveref}
\usepackage{multicol}

%--------Theorem Environments--------
%theoremstyle{plain} --- default
\newmdtheoremenv{theorem}{Theorem}[subsection]
\newmdtheoremenv{corollary}[theorem]{Corollary}
\newtheorem{proposition}[theorem]{Proposition}
\newmdtheoremenv{lemma}[theorem]{Lemma}
\newtheorem{conj}[theorem]{Conjecture}
\newtheorem{quest}[theorem]{Question}

\theoremstyle{definition}
\newmdtheoremenv{definition}[theorem]{Definition}
\newmdtheoremenv{definitions}[theorem]{Definitions}
\newtheorem{example}[theorem]{Example}
\newtheorem{algorithm}[theorem]{Algorithm}
\newtheorem{rmk}[theorem]{Remark}
\newtheorem{exercise}[theorem]{Exercise}

%----------- my commands -------------
\newcommand{\C}{\mathbb{C}}
\newcommand{\R}{\mathbb{R}}
\newcommand{\Q}{\mathbb{Q}}
\newcommand{\Z}{\mathbb{Z}}
\newcommand{\N}{\mathbb{N}}
\renewcommand{\P}{\mathbb{P}}
\newcommand{\contr}{\Rightarrow \Leftarrow}
\newcommand{\HRule}[1]{\rule{\linewidth}{#1}}
\newcommand{\s}{\(\,\)}
\renewcommand{\t}[1]{\text{#1}}
\newcommand{\ang}[1]{\left\langle #1 \right\rangle}

% Meta data
\setstretch{1.2}

\begin{document}

\pagestyle{fancy}
\fancyhf{}
\fancyhead[LOE]{\slshape{\textbf{Vincent Ferrara}}}
\fancyhead[ROE]{\thepage}

\section*{HW 1}
\subsection*{Problem 1}
\subsection*{1.}
\[\{(1,1),(1,-1)\}\]
\subsection*{2.}
\[\{(1,-1,1),(1,-1,-1),
   (-1,1,1),(-1,1,-1)
\}\]
\subsection*{3.}
\[A_n=\{(a_1,a_2,\dots,a_n,\dots) \in \Omega | \sum\limits_{i=1}^na_i=0\}\]
\subsection*{4.}
\[\bigcup\limits_{n=1}^\infty A_n\]
\subsection*{5.}
No for example the infinite sequence of +1'S

\newpage
\subsection*{Problem 2}
\subsection*{1.}
Since $\Omega=\bigcup\limits_{n \in N}\{n\}$ by definition 
\begin{align*}
    \P(\Omega)=
    \sum\limits_{n \in N}\P(\{n\})&=1\\
    \sum\limits_{n \in N}C_sn^{-s}&=1\\
    C_s\sum\limits_{n \in N}n^{-s}&=1\\
    C_s&=\dfrac{1}{\sum\limits_{n \in N}n^{-s}}
\end{align*}
\subsection*{2.}
Since  $A=\{2,4,6,\dots\}=\bigcup\limits_{n \in N}\{2n\}$ by definition 
\begin{align*}
    \P(A)
    &=\sum\limits_{n \in N}\P(\{2n\})\\
    &=\sum\limits_{n \in N}C_s(2n)^{-s}\\
    &=C_s2^{-s}\sum\limits_{n \in N}n^{-s}\\ \intertext{From 1. we know $C_s=\dfrac{1}{\sum\limits_{n \in N}n^{-s}}$} 
    &=\dfrac{1}{2^{s}}
\end{align*}
\subsection*{3.}
Since  $A=\{q,2q,3q,4q,\dots\}=\bigcup\limits_{n \in N}\{qn\}$ by definition
\begin{align*}
    \P(A)
    &=\sum\limits_{n \in N}\P(\{pn\})\\
    &=\sum\limits_{n \in N}C_s(pn)^{-s}\\
    &=C_sp^{-s}\sum\limits_{n \in N}n^{-s}\\ \intertext{From 1. we know $C_s=\dfrac{1}{\sum\limits_{n \in N}n^{-s}}$} 
    &=\dfrac{1}{p^{s}}
\end{align*}

Do the same thing for $B$ and we get $\P(B)=\dfrac{1}{q^s}$

Do the same thing for $A\cap B = \{pq,2pq,3pq,\dots\}$ and we get $\P(A\cap B)=\dfrac{1}{(pq)^s}$

\[\P(A)\P(B)=\P(A\cap B)=\dfrac{1}{(pq)^s}\]
Thus by definition event $A$ and event $B$ are independent with respect to this measure $\P$

\newpage
\subsection*{Problem 3}
\subsection*{1.}
First consider $A \cap B=(A^c \cup B^c)^c$
\begin{align*}
    \P(A\cap B)&=\P((A^c \cup B^c)^c)\\
               &=1-\P(A^c \cup B^c)\\ \intertext{Proven in class $\P(A^c \cup B^c) \leq \P(A^c)+\P(B^c)$}
               &\geq 1-\P(A^c)-\P(B^c)\\
               &= 1-1+\P(A)-1+\P(B)\\
               &=\P(A)+\P(B)-1
\end{align*}
Thus $\P(A\cap B) \geq \P(A)+\P(B)-1$\\
\phantom{text}\\
Now consider $(A\cap B)\cap C$ apply the proof above to these sets.
\begin{align*}
    \P((A\cap B)\cap C)&\geq \P(A\cap B) + \P(C)-1\\ \intertext{Now apply proof again}
    &\geq \P(A)+\P(B)\P(C)-2
\end{align*}
Therefore $\P((A\cap B)\cap C)\geq\P(A)+\P(B)\P(C)-2$
\subsection*{2.}

\newpage
\subsection*{Problem 4}
\begin{align*}
    B_1&=\{1,2,3,4,5,6\}\\
    B_2&=\{2,4,6\}\\
    B_3&=\{3,6\}
\end{align*}
\subsection*{1.}
\[
\P(6)=\P(6 \mid B_1)P(B_1)+\P(6 \mid B_2)P(B_2)+\P(6 \mid B_3)P(B_3)
     =\dfrac{1}{3}\left( \dfrac{1}{6}+\dfrac{1}{3}+\dfrac{1}{2} \right) = \dfrac{1}{3}
\]
\subsection*{2.}
\[
\P(B_2 \mid 6)=\dfrac{\P(B_2 \cap 6)}{\P(6)}=\dfrac{\P(6 \mid B_2)\P(B_2)}{\frac{1}{3}}=\dfrac{\frac{1}{3}\times\frac{1}{3}}{\frac{1}{3}}=\dfrac{1}{3}
\]
\subsection*{3.}
\[
\P(6 \text{ on second} \mid 6 \text{ on first}) = \dfrac{\P(6 \text{ on second} \cap 6 \text{ on first})}{\P(6)}=\dfrac{\frac{1}{3}(\P(6 \mid B_1)^2+\P(6 \mid B_2)^2+\P(6 \mid B_3)^2)}{\frac{1}{3}}
\]
\[
=\left( \dfrac{1}{36}+\dfrac{1}{9}+\dfrac{1}{4} \right)= \dfrac{7}{18}
\]
\subsection*{4.}
\[
\P(B_2 \mid \text{diff})=\dfrac{\P(B_2 \cap \text{diff})}{\P(\text{diff})}=\dfrac{\P(\text{diff}\mid B_2)\P(B_2)}{\P(\text{diff})}=\dfrac{\frac{2}{3}\times\frac{1}{3}}{\frac{1}{3}(\P(diff \mid B_1)+\dots)}=\dfrac{\frac{2}{3}}{(\frac{5}{6}+\frac{2}{3}+\frac{1}{2})}=\dfrac{1}{3}
\]
\subsection*{5.}
\begin{align*}
\P(\text{Same from two diff bags})&=\\
&=\P(\text{Same from two diff} \mid B_1,B_2)\P(B_1,B_2)\\
&+\P(\text{Same from two diff} \mid B_1,B_3)\P(B_1,B_3)\\
&+\P(\text{Same from two diff} \mid B_2,B_3)\P(B_2,B_3)\\
\\
\P(\text{Same from two diff} \mid B_1,B_2)&=\frac{1}{6}\\
\P(\text{Same from two diff} \mid B_1,B_3)&=\frac{1}{6}\\
\P(\text{Same from two diff} \mid B_2,B_3)&=\frac{1}{6}\\
\\
&=\dfrac{1}{3}\times\dfrac{1}{2}=\dfrac{1}{6}
\end{align*}
\subsection*{6.}
\begin{align*}
\P(B_1B_2 \mid \text{Same from two diff})=\dfrac{\P(\text{Same from two diff} \mid B_1,B_2)\P(B_1B_2)}{\P(\text{Same from two diff bags})}=\dfrac{1}{3}
\end{align*}
\end{document}